\documentclass[11pt, a4paper]{report}

\usepackage[utf8]{inputenc}
\usepackage[T1]{fontenc}
\usepackage[french]{babel}

% Appel des packages dans le dossier preambule
\def\packagepath{../../../../preambule}   % path du package princial

\usepackage{\packagepath/page_layout} % <--- Nouveau
\usepackage{\packagepath/config_info}
\usepackage{\packagepath/cadres_cours}
\usepackage{\packagepath/zones_reponse}
\usepackage{\packagepath/config_ds}
\usepackage{\packagepath/config_maths}

% --- PILOTAGE ---
\def\visibleornot{invisible}    % visible or invisible


\def\documentpath{./Sources_Latex}
\graphicspath{{./Images/}}


\def\ptsexoA{0}

\begin{document}

%%%%%%%%%%%%%%%%%%%%%%%%%%%
\def\level{BTS}  % Classe à droite
\def\course{CIEL 1}
\def\chaptername{Intégration} % Titre au centre
\def\docname{C107~--~TP4}        % Identifiant à gauche

\pagestyle{DS_FP}
%%%%%%%%%%%%%%%%%%%%%%%%%%%%%%%%%

\NomPrenomNote{}

%%=============================================================
\section*{PARTIE A\@: Détermination expérimentale du paramètre}

On cherche à modéliser un signal par la fonction suivante:  

\[f(x) =  p(x) = kx^2 + x + m\ln(x) \qquad \text{ (où $k>0$ et $m$ sont les paramètres à déterminer).}\]

Le signal doit respecter les contraintes physiques suivantes:

\begin{itemize}
    \item la courbe doit passer par le point $A(1;2)$.
    \item le signal doit présenter un maximum à l'instant $x=1$.
\end{itemize}

\begin{enumerate}
    \item Tracer la courbe de $f$ pour différentes valeurs de $k$ et de $m$. 
    \item Utiliser les contraintes pour déterminer les valeurs de $k$ et $m$. 
    \begin{blocvisible}[height=2]
        En faisant varier les valeurs de $k $ et $m$, on peut observer que la courbe de $f$ passe par le point $A(1;2)$ lorsque $k=1$ et que le signal présente un maximum à l'instant $x=1$ lorsque $m=-3$.
    \end{blocvisible}
\end{enumerate}



\section*{PARTIE B\@: Étude du signal identifié}

Utilisez maintenant les valeurs de $k$ et $m$ déterminées précédemment pour étudier le signal modélisé par la fonction $f(x) = x^2 + x - 3\ln(x)$.

\begin{enumerate}[resume]
    \item Tracer la courbe de $f$ et observer son comportement. 
    \item Affichez la fonction dérivée $f'$  
    \item Déterminer le signe de $f'$ et en déduire les variations de $f$ sur $\oo{0}{+\infty}$. 
    \begin{blocvisible}[height=5]
        \begin{minipage}{0.58\linewidth}
            \[f'(x) = 2x + 1 - \dfrac{3}{x} = \dfrac{2x^2 + x - 3}{x} = \dfrac{(2x+3)(x-1)}{x}\] 
        
            qui s'annule et change de signe en $x=1$.

            La fonction $f$ est décroissante sur l'intervalle $\oo{0}{1}$ et croissante sur l'intervalle $\oo{1}{+\infty}$.
        \end{minipage}\hfill
        \begin{minipage}{0.4\linewidth}
            \begin{tikzpicture}
   % lgt: largeur 1ère col, espcl: espace entre les colonnes de valeurs
   \tkzTabInit[lgt=2, espcl=1.8]{$x$ / 1 ,  $f'(x)$ / 1, $f$ / 1.5}
      {0, 1, $+\infty$} % Valeurs de x
      
   % Ligne du signe de (2x+3) : toujours positif sur cet intervalle
   
   % Ligne du signe de (x-1) : s'annule en 1
   
   % Ligne du signe de x : toujours positif
   
   % Signe de la dérivée (produit/quotient des lignes précédentes)
   \tkzTabLine{ d, -, z, +, }
   
   % Variations de f
   \tkzTabVar{+ / $+\infty$, - / $2$, + / $+\infty$}
\end{tikzpicture}
\end{minipage}
    \end{blocvisible}
        
    
    \item Placer un point $N\left(x_N; f(x_N)\right)$ sur la courbe et observer la valeur de $f(x_N)$ pour différentes positions de $N$.        
    \item Déplacez le point $N$ vers la droite (grandes valeurs de $x$) et observez la tendance de $f(N)$. 
    \item Conjecturez la limite de $f(x)$ lorsque $x$ tend vers $+\infty$.  
    
    \begin{blocvisible}[height=3]
        En déplaçant le point $N$ vers la droite, on observe que la valeur de $f(x)$ augmente de plus en plus. En conjecturant la limite de $f(x)$ lorsque $x$ tend vers $+\infty$, on peut supposer que $\lim_{x \to +\infty} f(x) = +\infty$.

        Pour justifier cette conjecture, on peut calculer la limite de $f(x)$ lorsque $x$ tend vers $+\infty$:
        \[\lim_{x \to +\infty} f(x) = \lim_{x \to +\infty} (x^2 + x - 3\ln(x))=+\infty\]
    \end{blocvisible}
\end{enumerate}


\pagebreak
\thispagestyle{DS_LP}

\section*{PARTIE C\@: Intégration et Valeur Moyenne}

On souhaite évaluer la tension moyenne délivrée par le composant sur un intervalle $\ff{a}{b}$, avec $0<a<b$.

\begin{enumerate}[resume]
    \item Créer les curseurs $a$ et $b$ avec les valeurs initiales $a=0,1$ et $b=6$.  
    \item Déterminer graphiquement l'aire sous la courbe de $f$ entre $a$ et $b$ avec la commande 
    
    \verb|Aire = Intégrale(f, a, b)|.  
    \begin{blocvisible}[height=2]
        D'après Géogébra, l'aire sous la courbe de $f$ entre $a=0,1$ et $b=6$ est d'environ $1,25$ unités d'aires.
    \end{blocvisible}
        
    \item Observer comment l'aire change lorsque vous modifiez les valeurs de $a$ et $b$. 
    \begin{blocvisible}[height=3]
        Lorsque vous augmentez $b$, l'aire sous la courbe de $f$ entre $a$ et $b$ augmente, car vous incluez plus de la courbe. Inversement, lorsque vous augmentez $a$, l'aire diminue, car vous excluez une partie de la courbe.
    \end{blocvisible}

    \item A l'aide de Géogébra, déterminer m'expression d'une primitives $F$ de $f$. 
    \begin{blocvisible}[height=2]
        \[F(x) = \frac{1}{3}\cdot x^3 + \frac{1}{2} \cdot x^2 - 3x\cdot \ln(x) + 3x\]
    \end{blocvisible}

    \item Calculer l'aire sous la courbe de $f$ entre $a=0,1$ et $b=6$ à l'aide de l'intégrale. 
    \begin{blocvisible}[height=3]
        \[\int_{0,1}^6 f(x) \dx = F(6) - F(0,1) \approx 1,25\]

        L'aire calculée à l'aide de l'intégrale est d'environ $1,25$ unités d'aires, ce qui correspond à l'aire trouvée graphiquement.
    \end{blocvisible}

    \item Déterminer la valeur moyenne $\mu$ de la fonction $f$ sur l'intervalle $\ff{a}{b}$.
    \begin{blocvisible}[height=3]
        \[\mu = \frac{1}{b-a}\cdot \int_a^b f(x) \dx = \frac{1}{6-0,1} \cdot 1,25 = 0,21\]

        La valeur moyenne de la fonction $f$ sur l'intervalle $[0,1 ; 6]$ est d'environ $0,21$ volts.
    \end{blocvisible}

    \item Créer un rectangle d'équivalence avec les points \verb|(a,0)|, \verb|(b,0)|,\verb|(b, moy)| et \verb|(a, moy)|. 
    
    Utilisez l'outil \verb|Polygone| pour le colorer. 

    \item En conservant $a=0$, déterminer la valeur de $b$ pour laquelle la valeur moyenne $\mu$ de $f$ sur l'intervalle $\ff{a}{b}$ est égale à $0,5$  (à $10^{-2}$ près). 
    \begin{blocvisible}[height=3]
        A l'aide de Géogébra, on fixera $a=0$ et on fera varier $b$ jusqu'à ce que la valeur moyenne $\mu$ soit égale à $0,5$ volts.

        On obtient une valeur de $b$ d'environ $2,5$.
    \end{blocvisible}
\end{enumerate}


\end{document}
